\chapter[Conclusão]{Conclusão}
\label{cap:conclusao}

Este trabalho teve como objetivo desenvolver o RECLin-PT, modelo de extração de
relações em notas clínicas em português sobre o corpus SemClinBr, com ênfase na
relação \texttt{negation\_of}. A resposta à pergunta de pesquisa, sobre como
desenvolver esse modelo, tem três componentes. O primeiro é a cadeia de dados: a
reconstrução do corpus a partir dos XML com o texto alinhado às anotações, as
partições congeladas e o gerador de pares candidatos, que tornam a tarefa definida e
comparável. O segundo é o modelo: a representação do par por marcadores de entidade e
a arquitetura Pair-Aware linear, que lê explicitamente as representações do par e é
ajustada junto com o \emph{encoder}. O terceiro é o procedimento de definição da
configuração: a comparação de \emph{encoders} e de estratégias de modelagem, sob os
mesmos dados e o mesmo protocolo, como base para escolher a configuração do modelo.

Os quatro primeiros objetivos específicos foram cumpridos: os dados foram preparados,
as relações foram caracterizadas e os pares candidatos gerados, o RECLin-PT foi
desenvolvido com a Pair-Aware linear e comparado com estratégias alternativas, e as
configurações foram avaliadas no DEV com três \emph{encoders}. O quinto foi cumprido
em parte: a configuração de referência atual é a Pair-Aware linear com o
BioBERTpt-all e a semente $42$, e o congelamento e a avaliação única no teste estão
previstos no protocolo e na emenda proposta, mas o critério de seleção da
configuração final continua em aberto, e essa configuração, inclusive o
\emph{encoder}, ainda não foi selecionada, congelada nem avaliada no teste. Os
resultados que orientaram essas decisões são provisórios e delimitados ao DEV e a uma
semente: na classe-alvo, nem a troca de \emph{encoder} nem a troca entre a Pair-Aware
linear e o Two-Stage produziram diferença detectável; a estratégia influenciou
sobretudo as outras classes, e o \emph{encoder}, sobretudo \texttt{associated\_with}.

\section{Principais resultados}

Nos três \emph{encoders}, a Pair-Aware linear tem F1 de \texttt{negation\_of} maior
que o do baseline ($0{,}829$, $0{,}826$ e $0{,}818$ com o BioBERTpt-all, o
BioBERTpt-clin e o BERTimbau-base, contra $0{,}802$, $0{,}814$ e $0{,}800$), com
diferença detectada só com o BioBERTpt-all, e macro-F1 e F1 de
\texttt{associated\_with} maiores, com diferença detectada nos três, inclusive quando
se reamostram documentos inteiros. O Two-Stage alcança F1 de \texttt{negation\_of}
equivalente ($0{,}825$ a $0{,}832$), sem diferença detectada para a Pair-Aware
linear, mas depende de um limiar escolhido no DEV, que, com os dois BioBERTpt, ficou
no extremo da grade e sacrificou a classe \texttt{associated\_with}. Entre
\emph{encoders}, nenhuma das nove comparações no F1 de \texttt{negation\_of} detectou
diferença. O BioBERTpt-clin teve F1 de \texttt{associated\_with} maior na Pair-Aware
linear e no baseline, e o BERTimbau-base, com cerca de 61\% dos parâmetros, ficou
equivalente na classe-alvo. A análise dos falsos positivos de
\texttt{associated\_with} mostrou que cerca de um quarto deles são pares cuja
direção inversa está anotada. Para o desenvolvimento do RECLin-PT, esses resultados
são compatíveis com a Pair-Aware linear como arquitetura, sem a demonstrar superior
na classe-alvo, e não justificam trocar o BioBERTpt-all, o \emph{encoder} do
protocolo, que segue como referência.

\section{Limitações}

Essas conclusões não podem ser tomadas como definitivas. Do ponto de vista
experimental, todos os números são de DEV, o mesmo conjunto usado nas escolhas de
época e de limiar; há uma única semente; a designação da Pair-Aware linear como
arquitetura oficial não é confirmatória; as variantes de \emph{encoder} estão fora do
pré-registro; e a configuração final do RECLin-PT não foi congelada nem avaliada no
teste. Se a emenda proposta for aprovada, o teste avaliará só o modelo final, de modo
que as diferenças entre estratégias e \emph{encoders} continuarão apoiadas apenas no
DEV. Somam-se a isso limitações arquiteturais (três classes e modelos de porte
\emph{base}), dos dados (inconsistências de anotação e uma única partição), de escopo
(entidades previamente anotadas, sem reconhecimento automático) e de reprodutibilidade
(pontos de procedência das execuções). Essas categorias são detalhadas na
Seção~\ref{sec:limitacoes}, e as ameaças à validade, na seção seguinte a ela.

\section{Contribuições}

As principais contribuições são: (i) o desenvolvimento do RECLin-PT, com a
Pair-Aware linear como arquitetura oficial e a comparação, sob o mesmo protocolo e
com três \emph{encoders}, de estratégias de modelagem para a extração de relações
clínicas em português, do classificador direto a uma arquitetura em dois estágios;
(ii) a comparação, sob configuração idêntica e em três estratégias, entre
\emph{encoders} clínicos e um geral na extração de relações sobre o SemClinBr, que
subsidia a escolha do \emph{encoder} do modelo e oferece evidência inicial sobre a
lacuna descrita no Capítulo~\ref{cap:relacionados}; (iii) a cadeia de dados que
sustenta essas comparações, com os dados reconstruídos dos XML e o texto alinhado aos
offsets de anotação, partições em nível de documento congeladas por \emph{hashes} e um
gerador determinístico de pares candidatos compartilhado por todas as estratégias,
com o teto de \emph{recall} quantificado; (iv) uma caracterização do corpus que
evidencia o desbalanceamento das relações e a quase adjacência dos pares de negação;
e (v) um protocolo pré-registrado e uma matriz de rastreabilidade que ligam cada
escolha e cada número do texto ao artefato de origem, complementados por diagnósticos
reproduzíveis dos dados e das execuções: a sensibilidade da partição à semente, a
detecção de documentos quase duplicados, controles informados pelas anotações, que
medem o atalho oferecido pelo tipo semântico \texttt{Negation}, e a auditoria das
execuções, que recalcula as métricas a partir das predições gravadas e confere os
vínculos entre elas. No contexto do campus, o trabalho dá continuidade às pesquisas
em reconhecimento de entidades clínicas de \citeonline{torres_2024_ner} e
\citeonline{ramos_2025_clinpthyb}, que apontavam a extração de relações como
possibilidade de continuidade.

\section{Trabalhos futuros}

Os trabalhos futuros dividem-se em dois grupos. O primeiro reúne as etapas
necessárias para concluir a avaliação do RECLin-PT, que decorrem diretamente das
limitações experimentais e de reprodutibilidade:

\begin{itemize}
  \item aprovar e registrar, antes de qualquer avaliação no teste, a emenda ao
        protocolo redigida como proposta (E3), depois de definidos os critérios
        estatísticos das comparações e o critério de seleção da configuração final,
        ambos ainda em aberto;
  \item conduzir os experimentos de aprimoramento só com treino e DEV, registrando
        todas as tentativas, e comparar cada uma com a configuração de referência e
        com os baselines do mesmo \emph{encoder};
  \item selecionar e congelar a configuração final, conferir a proveniência do modelo
        final, inclusive o \emph{checkpoint}, e registrar a identidade do código;
  \item avaliar no teste, uma única vez, apenas o modelo final congelado, e relatar
        essa avaliação também sem as três notas do teste que têm cópia no treino, como
        análise secundária pré-registrada;
  \item repetir os experimentos com mais sementes, para separar o efeito das
        estratégias e dos \emph{encoders} da variação do treino;
  \item calcular os controles do protocolo (baseline calibrado e filtro léxico) também
        para o BioBERTpt-clin e o BERTimbau-base.
\end{itemize}

O segundo grupo reúne possibilidades de evolução do modelo e de pesquisa futura:

\begin{itemize}
  \item integrar NER e RE em um fluxo completo, com uma etapa automática de
        reconhecimento de entidades, como as investigadas por
        \citeonline{torres_2024_ner} e \citeonline{ramos_2025_clinpthyb}, à frente do
        RECLin-PT, e avaliar a extração de relações nesse cenário fim a fim, em que
        os erros de reconhecimento se propagam para as relações;
  \item investigar outras relações clínicas, além das duas anotadas no SemClinBr,
        o que exigiria outros corpora ou anotação adicional, e examinar, para
        \texttt{associated\_with}, o efeito da direção anotada, por exemplo avaliando
        a classe sem direção, em protocolo próprio;
  \item aprimorar a identificação das relações de negação, analisando os erros por
        tipo de gatilho e por distância entre os argumentos e realizando a adjudicação
        cega dos falsos positivos cuja pista não tem relação anotada, para separar
        erros do modelo de possíveis omissões da anotação;
  \item investigar estratégias adicionais de modelagem e de representação: executar
        o \emph{fine-tuning} restrito previsto no protocolo; testar, em novo
        protocolo, critérios de limiar que considerem as três classes, grades que não
        terminem no valor escolhido e orçamentos de treino maiores, já que a melhor
        época coincidiu com o teto; e avaliar, em uma variante experimental separada e
        com protocolo próprio, tipos semânticos agregados na entrada, com o tipo
        \texttt{Negation} mascarado, já que esse tipo, sozinho, recupera quase todas
        as pistas (Seção~\ref{sec:controles_tipo}).
\end{itemize}
